% 谷神星（综述）
% license CCBYSA3
% type Wiki

本文根据 CC-BY-SA 协议转载翻译自维基百科\href{https://en.wikipedia.org/wiki/Ceres_(dwarf_planet)}{相关文章}。

\begin{figure}[ht]
\centering
\includegraphics[width=6cm]{./figures/ed4afc93bd247d5b.png}
\caption{Dawn 探测器于 2015 年 5 月拍摄的 Ceres 图像。两个明亮斑点分别为 Haulani 撞击坑（右）和 Oxo 撞击坑坑底（左）。} \label{fig_Ceres_1}
\end{figure}
Ceres（小行星编号：1 Ceres）是一颗位于 Mars 与 Jupiter 轨道之间主小行星带中的矮行星。它是人类发现的首颗小行星，于 1801 年 1 月 1 日由 Giuseppe Piazzi 在 Sicily 的 Palermo Astronomical Observatory 发现，并在当时被宣布为一颗新行星。此后，Ceres 被重新分类为小行星，近年又被认定为矮行星——它是唯一一颗轨道未超出 Neptune 之外的矮行星，也是所有不具备卫星的矮行星中最大的一颗。

Ceres 的直径约为 Moon 的四分之一。由于其体积很小，即使在最亮时，除非在极其黑暗的天空条件下，也无法用肉眼看到。其视星等范围为 6.7 至 9.3，并在每隔 15 至 16 个月发生一次会合（即靠近 Earth）时达到峰值。因此，即使使用最强大的望远镜，其表面特征也只能勉强识别，在 NASA Dawn 无人探测器于 2015 年抵达并进入 Ceres 轨道任务之前，人们对其所知甚少。

Dawn 的观测发现，Ceres 的表面由水、冰以及碳酸盐和黏土等含水矿物混合组成。重力数据表明，Ceres 的内部可能呈部分分化结构，拥有一个泥状（冰—岩）地幔/核心，以及一层密度较低但更坚固的地壳（其体积最多约 30\% 为冰）。虽然 Ceres 可能缺乏内部液态水海洋，但盐水仍在外地幔中流动并抵达表面，从而大约每五千万年形成一次如 Ahuna Mons 这样的低温火山。这使 Ceres 成为距离 Sun 最近的已知低温火山活跃天体。Ceres 具有极其稀薄且短暂的水蒸气大气，由其表面局部区域喷发形成。
\subsection{历史}  
\subsubsection{发现}  
在 18 世纪日心说被接受到 1846 年 Neptune 被发现之间的多年里，一些天文学家认为数学定律预测在 Mars 和 Jupiter 轨道之间存在一颗隐藏或失踪的行星。1596 年，理论天文学家 Johannes Kepler 认为，只有在添加两颗行星——一颗位于 Jupiter 与 Mars 之间，另一颗位于 Venus 与 Mercury 之间——时，行星轨道之间的比例才会符合“上帝的设计”\(^{[19]}\)。其他理论家，例如 Immanuel Kant，思考该空隙是否由 Jupiter 的引力造成；1761 年，天文学家兼数学家 Johann Heinrich Lambert 问道：“谁又能知道是否已有行星从 Mars 与 Jupiter 之间的广阔空间中消失？天体是否如同 Earth 一样，强者磨损弱者，Jupiter 与 Saturn 是否注定永远掠夺它们？”\(^{[19]}\)

1772 年，德国天文学家 Johann Elert Bode 引用 Johann Daniel Titius 的研究，发表了一条后来被称为 Titius–Bode 定律的公式，该公式似乎能预测已知行星的轨道，但在 Mars 与 Jupiter 之间出现无法解释的空隙\(^{[19][20]}\)。该公式预测 Sun 附近约 2.8 AU（约 4.2 亿 km）处应存在另一颗行星\(^{[20]}\)。随着 William Herschel 于 1781 年在预测距离附近发现 Saturn 之外的 Uranus，这一定律获得更多信任\(^{[19]}\)。1800 年，由德国天文期刊《Monthly Correspondence》主编 Franz Xaver von Zach 领导的一组人向 24 名经验丰富的天文学家发出请求，将他们称为“天上警察”\(^{[20]}\)，希望他们共同努力，开始系统搜索这颗预期行星\(^{[20]}\)。尽管他们未发现 Ceres，但后来发现了小行星 Pallas、Juno 和 Vesta\(^{[20]}\)。

受邀参与该搜索计划的天文学家之一是 Giuseppe Piazzi，他是 Sicily Palermo 学院的一名天主教神父。在收到邀请函之前，Piazzi 于 1801 年 1 月 1 日发现了 Ceres\(^{[21]}\)。他当时在寻找“la Caille 氏黄道星表中的第 87 颗恒星”，却发现“前面还有另一颗”\(^{[19]}\)。Piazzi 并非找到了一颗恒星，而是一颗运动中的类星体，他起初认为那是一颗彗星\(^{[22]}\)。Piazzi 共观测 Ceres 24 次，最后一次观测发生在 1801 年 2 月 11 日，后因生病中断。他于 1801 年 1 月 24 日通过信件向两位天文学家宣布了他的发现：一位是来自 Milan 的同胞 Barnaba Oriani，另一位是 Berlin 的 Bode\(^{[23]}\)。他将其报告为彗星，但表示“由于它的运动如此缓慢且近乎均匀，我多次想到它可能比彗星更重要”\(^{[19]}\)。4 月，Piazzi 将完整观测结果送给 Oriani、Bode 和法国天文学家 Jérôme Lalande。这些信息发表于 1801 年 9 月的《Monatliche Correspondenz》\(^{[22]}\)。

当时，Ceres 的视位置已经发生改变（主要由于 Earth 绕 Sun 运动），且过于靠近 Sun 的眩光，其他天文学家无法确认 Piazzi 的观测。到了年底，Ceres 应该能再次被观测到，但由于已经过去很长时间，其位置难以预测。为重新定位 Ceres，24 岁的数学家 Carl Friedrich Gauss 开发了一套高效的轨道确定方法\(^{[22]}\)。他在几周内预测出 Ceres 的轨迹，并将结果发送给 von Zach。1801 年 12 月 31 日，von Zach 与同为“天上警察”的 Heinrich W. M. Olbers 在预测位置附近再次发现 Ceres，并继续记录其位置\(^{[22]}\)。Ceres 与 Sun 相距 2.8 AU，几乎完美符合 Titius–Bode 定律；但当 1846 年 Neptune 被发现时，其实际位置比预测近 8 AU，多数天文学家认为该定律仅是巧合\(^{[24]}\)。

早期观测者只能大致估算 Ceres 的大小。Herschel 在 1802 年低估其直径为 260 km（160 mi）；1811 年，德国天文学家 Johann Hieronymus Schröter 将其高估为 2,613 km（1,624 mi）\(^{[25]}\)。20 世纪 70 年代，红外光度测量使其反照率测定更为精确，Ceres 的直径得出与真实值 939 km（583 mi）相差不到 10\% 的结果\(^{[25]}\)。















